\documentclass[11pt]{article}
\usepackage[margin=1in]{geometry}
\usepackage{titlesec}
\usepackage{url}
\usepackage{hyperref}
\usepackage{amsmath}
\usepackage{enumitem} % in preamble
\setlist[itemize]{leftmargin=*, nosep, topsep=0pt, partopsep=0pt}

\newcommand{\memoheader}[4]{
  \noindent
  \textbf{From:} #2\\ 
  \textbf{To:} #1\\
  \textbf{Date:} #3\\
  \textbf{Subject:} #4\\[1ex]\hrule\vspace{1.5ex}
}

\begin{document}
\memoheader{Prof. Barbara Linke}
            {Miguel Cuaycong}
            {\today}
            {EME 150B: Gear Fundamentals}

\section{Nomenclature}
\begin{itemize}
    \item Pitch circle 
    \item Number of teeth $\equiv$ N
    \item Pitch diameter $\equiv$ diameter of pitch circle
    \item Pinion $\equiv$ smaller of the two mating gears
    \item circular pitch $\equiv$ p; arc length from two adjacent tooth
    \item Module $\equiv m = \frac{\text{Pitch Diameter}}{N}$ ; ratio of pitch diameter to number of teeth.
    \item Diametral pitch P $\equiv \frac{1}{m}$
    \item Addendum $\equiv$ a; radial distance between the top land and the pitch circle
    \item Dedendum $\equiv$ b; distance between bottom land and pitch circle
\end{itemize}
\section{No slip assumption}
Pitch circles roll into one another without slipping
\\
Remember:
\begin{itemize}
    \item $\omega$ [rad/s]
\end{itemize}
\subsection{pitch line velocity $V$}
\begin{align*}
    V_1 = V_2 &\rightarrow |r_1\omega_1| = |r_2\omega_2|
    \\
    |\frac{\omega_1}{\omega_2}| &= \frac{r_1}{r_2}
\end{align*}
\subsection{Torque}
$$T = \text{cross product}(r, F) [\text{Nm}]$$
Torque ratio: $\frac{T_1}{T_2}$; The convention is for the pinion (driving gear) to be gear ?1.
\subsection{Power}
$$H = T\omega \rightarrow H(t)=FV$$
Power ratio: $\frac{H_1}{H_2}$
\subsection{Involute properties}
Insert hand-drawn sketches
\begin{align*}
    r_{b,i}&=r_i \cos \phi
    \\
    p_b &= p_c \cos \phi
\end{align*}
Where: $\phi \equiv$ pressure angle
\section{Interference}
\begin{align*}
    N_p=\frac{2k}{3 \sin^2 \phi}\bigg( 1 + \sqrt{1 + 3 \sin^2 \phi} \bigg)
\end{align*}
\section{Gear train}
let $n = \omega = \frac{rev}{time}$
\\
Two gears mated have $n$'s that are proportional to each other e.g:
\\
$n_1 = An_2$
\begin{align*}
    V_1&=V_2
    \\
    r_1 \omega_1 &= r_2 \omega_2
    \\
    \omega_1 &= \frac{r_2}{r_1} \omega_2
    \\
    \omega_1 &= \frac{d_2}{d_1} \omega_2
    \\
    \omega_1= n_1 &= \frac{N_2}{N_1} n_2; d = \frac{N}{P} \text{ where by rule 2: P is the same for all mated gears.}  
\end{align*}
This mathematical relationship will be the basis of gear train with k amount of gears. the kth gear's radial speed of $n_k$ can be expressed relative to the driving gear $n_1$ as:
\\~\\
\begin{align*}
    n_k = \frac{N_{k-1}}{N_{k}} \frac{N_{k-2}}{N_{k-1}} \dots \frac{N_1}{N_2}n_1
\end{align*}

\section*{Notes}
\begin{itemize}
    \item the overall train value is the product of the gear set train values.
    \item if one gear set achieves up to $e_{set}=\frac{1}{10}$, then two gear sets achieve $e_{set,1} e_{set,2}$
    \item Can use $e_{set,1} e_{set,2}$ as a criteria: $e_{set,1} e_{set,2} < e_{desired}$
    \item Investigate Conjugate action relationship to the Involute curve
\end{itemize}

\end{document}